\section{Auswertung}
\label{sec:Auswertung}

Die Graphen werden sowohl mit Matplotlib \cite{matplotlib} als auch NumPy \cite{numpy} erstellt. Die Fehlerrechnung wird mithilfe von Uncertainties \cite{uncertainties} durchgeführt. Die Konstanten werden SciPy \cite{scipy} entnommen.

\subsection{Berechnung der theoretischen Werte}

Zu Beginn werden die theoretischen Landefaktoren der Energieniveaus aus den Übergängen der roten Cadmiumlinie ($^3D_2 \rightarrow^3P_1$) und der blauen Cadmiumlinie ($^3P_1 \rightarrow^3S_1$) berechnet. 
Zudem werden das Dispersionsgebiet $\Delta\lambda_D$ und das Auflösungsvermögen $A$ bestimmt.
Es werden dabei die Landefaktoren $g_im_i-g_jm_j$ aus Gleichung \ref{eq:gij} zu $g_{ij}$ zusammengefasst.
Für die rote Cadmiumlinie ergeben sich die Landefaktoren
\[
g_{i_.{norm}}=g_{j_.{norm}}=1\text{.}
\]
Daraus folgen für die Übergänge mit der Auswahlregel $\Delta m=0 \lor \pm 1$:
\begin{align*}
g_{ij_.{norm,\pi}}&=0\\
g_{ij_.{norm,\sigma}}&=\pm 1\\
\end{align*}
Für die blaue Cadmiumlinie ergeben sich die Landefaktoren
\begin{align*}
g_{i_.{anom}} &= 1,5\\
g_{j_.{anom}} &= 2 \text{.}
\end{align*}
Daraus folgen für die Übergänge:
\begin{align*}
g_{ij_.{anom,\pi}}&=\pm 0,5\\
g_{ij_.{anom,\sigma1}}&= \pm 1,5\\
g_{ij_.{anom,\sigma2}}&= \pm 2 \text{.}
\end{align*}
Der $\pi$-Übergang $m_2=0\rightarrow m_1=0$ ist trotz $\Delta m=0$ unterdrückt.
Da die $g_{ij_.{anom,\sigma}}$ sehr dicht beisammen liegen, werden bei der Messung nur Aufspaltungen mit $g_{ij_.{anom,\sigma}} = \pm 1,75$ gemessen.
Das Dispersionsgebiet ergibt sich für die rote Linie mit $\lambda_.{rot}=\SI{643.8}{\nano\metre}$ und für die blaue Linie mit $\lambda_.{blau}=\SI{480}{\nano\metre}$ bei einer Lumma-Gehrcke Platte mit den Dimensionen $L=\SI{120}{\milli\metre}$, $d=\SI{4}{\milli\metre}$ und den Brechungsindizes $n(\SI{480}{\nano\metre})=1,4635$ und $n(\SI{643.8}{\nano\metre})=1,4567$ und Formel \eqref{eq:disGebiet} zu:
\[
\Delta\lambda_D =
\begin{cases}
\SI{48.91}{\pico\metre} \text{, für }\lambda=\SI{643.8}{\nano\metre}\\
\SI{26.95}{\pico\metre} \text{, für }\lambda=\SI{480}{\nano\metre}
\end{cases}\text{.}
\]
Das Auflösungsvermögen berechnet sich nach Formel \eqref{eq:Aufloesung} zu:
\[
A = 
\begin{cases}
\num{2.09e5} \text{, für }\lambda=\SI{643.8}{\nano\metre}\\
\num{2.85e5} \text{, für }\lambda=\SI{480}{\nano\metre}
\end{cases}\text{.}
\]

\subsection{Kalibrierung des B-Feldes}

Es wird das B-Feld kalibriert. Hierzu wird die stärke des B-Feldes in Abhängigkeit von der Stromstärke I aufgenommen. Die Werte sind in Tabelle \ref{tab:B} aufgelistet und in Abbildung \ref{fig:B} graphisch dargestellt. Es wird eine lineare Ausgleichsrechnung der Form
\[
B(I) = aI+b
\]
durchgeführt. Dies liefert die Werte:
\begin{align*}
a &= \SI{60.24(28)}{\milli\tesla\per\ampere}\\
b &= \SI{0.1(3)}{\milli\tesla}
\end{align*}

\begin{figure}
	\centering
	\includegraphics[width=\linewidth-130pt,height=\textheight-130pt,keepaspectratio]{build/Bfeldkali.pdf}
	\caption{Die Kalibrationskurve des B-Feldes in Abhängigkeit der Stromstärke I.}
	\label{fig:B}
\end{figure}

\begin{table}
	\centering
	\caption{Die Messwerte für die B-Feldkalibrierung.}
	\label{tab:tabIB1}
	\sisetup{table-format=1.2}
	\begin{tabular}{S[table-format=2.1]S[table-format=1.3]S[table-format=2.1]S[table-format=1.3]}
		\toprule
		{$ I / \si{\ampere}$} & {$ B/ \si{\tesla}$} & {$ I / \si{\ampere}$} & {$ B/ \si{\tesla}$} \\
		\midrule
		0.0 & 0.000 & 11.0 & 0.661 \\
		1.0 & 0.062 & 11.5 & 0.685 \\
		2.0 & 0.122 & 12.0 & 0.730 \\
		3.0 & 0.192 & 12.5 & 0.762 \\
		4.0 & 0.230 & 13.0 & 0.786 \\
		5.0 & 0.300 & 13.5 & 0.813 \\
		6.0 & 0.348 & 14.0 & 0.846 \\
		7.0 & 0.425 & 14.5 & 0.880 \\
		8.0 & 0.480 & 15.0 & 0.906 \\
		9.0 & 0.550 & 16.0 & 0.961 \\
		10.0 & 0.600 & 17.0 & 1.014 \\
		\bottomrule
	\end{tabular}

	%\input{build/tabIB2.tex}
	\label{tab:B}
\end{table}

\subsection{Untersuchung des normalen Zeeman-Effekts}

Es wird der normale Zeeman-Effekt anhand der $\sigma$-Übergänge der roten Cadmiumlinie mit $\lambda=\SI{643.8}{\nano\metre}$ untersucht. Das Magnetfeld wird mit $\SI{10.5}{\ampere}$ induziert, was mit den Werten der Kalibrierung einem Magnetfeld von $\SI{633}{\milli\tesla}$ entspricht.\\
In Abbildung \ref{fig:rot_sigma} ist die Aufnahme der Spektrallinien mit und ohne Magnetfeld zu sehen. Die gemessenen Werte für $\Delta s$ und $\delta s$ sind zusammen mit den nach Formel \eqref{eq:deltaLambda} berechneten Werten für $\delta\lambda$ in Tabelle \ref{tab:rot_sigma} eingetragen. Die Abstände werden dabei digital aus den Bildern entnommen und sind deswegen in Pixel (px) angegeben. Zur Berechnung von $\delta\lambda$ wird der zuvor bestimmte Wert für $\Delta\lambda_D$ genutzt.
Über alle Werte von $\delta\lambda$ gemittelt ergibt sich:
\[
\delta\lambda_.{norm,\sigma} = \SI{11.95(11)}{\pico\metre}\text{.}
\]
Mit den Formeln \eqref{eq:gij} und \eqref{eq:EaufExp} ergibt sich für die $g_{ij}$:
\begin{equation}
g_.{ij}=\frac{h c}{\lambda^2 \cdot \mu_\text{B} B} \cdot \delta \lambda \label{eq:exgj}
\end{equation}
und damit
\[
g_{ij_.{norm,\sigma}} = \pm\num{0.976(10)}\text{.}
\]

\begin{figure}
	\centering
	\includegraphics[width=\linewidth-130pt,height=\textheight-130pt,keepaspectratio]{content/images/rot_sigma.jpg}
	\caption{Die Aufspaltung der roten Cadmium $\sigma$-Linien.}
	\label{fig:rot_sigma}
\end{figure}

\begin{table}
	\centering
	\caption{Die Messwerte für $\Delta s$ und $\delta s$, sowie die berechneten Werte $\delta\lambda$ für die Aufspaltung der roten Cadmium $\sigma$-Linien.}
	\label{tab:tabrotsigma}
	\sisetup{table-format=1.2}
	\begin{tabular}{S[table-format=3.0]S[table-format=3.0]S[table-format=1.2]}
		\toprule
		{$ \Delta s_{\text{norm},\sigma}/ px$} & {$ \delta s_{\text{norm},\sigma}/ px $} & {$ \delta\lambda_{\text{norm},\sigma}/ \si{\pico\meter}$} \\
		\midrule
		492 & 230 & 11.43 \\
		442 & 209 & 11.56 \\
		405 & 201 & 12.14 \\
		378 & 186 & 12.03 \\
		353 & 181 & 12.54 \\
		345 & 176 & 12.48 \\
		320 & 156 & 11.92 \\
		288 & 140 & 11.89 \\
		278 & 133 & 11.70 \\
		270 & 130 & 11.77 \\
		\bottomrule
	\end{tabular}

	\label{tab:rot_sigma}
\end{table}

\subsection{Untersuchung des anomalen Zeeman-Effekts}

Es wird der anomale Zeeman-Effekt anhand der $\sigma$- und $\pi$-Übergänge der blauen Cadmiumlinie mit $\lambda=\SI{480}{\nano\metre}$ untersucht. Zur Bestimmung von $\delta\lambda$ wird der passende Wert des zuvor bestimmten $\Delta\lambda_D$ genutzt. 

\subsubsection{Untersuchung der $\sigma$-Linien}

Das Magnetfeld wird mit $\SI{5.5}{\ampere}$ induziert, was mit den Werten der Kalibrierung einem Magnetfeld von $\SI{332}{\milli\tesla}$ entspricht. In Abbildung \ref{fig:blau_sigma} ist die Aufnahme der Spektrallinien mit und ohne Magnetfeld zu sehen. Die gemessenen Werte für $\Delta s$ und $\delta s$ sind zusammen mit den nach Formel \eqref{eq:deltaLambda} berechneten Werten für $\delta\lambda$ in Tabelle \ref{tab:blau_sigma} eingetragen.
Über alle Werte von $\delta\lambda$ gemittelt ergibt sich:
\[
\delta\lambda_.{anom,\sigma} = \SI{6.21(3)}{\pico\metre}\text{.}
\]
Mit Formel \eqref{eq:exgj} ergibt sich für die $g_{ij}$:
\[
g_{ij_.{anom,\sigma}} = \pm\num{1.714(13)}\text{.}
\]

\begin{figure}
	\centering
	\includegraphics[width=\linewidth-130pt,height=\textheight-130pt,keepaspectratio]{content/images/blau_sigma.jpg}
	\caption{Die Aufspaltung der blauen Cadmium $\sigma$-Linien.}
	\label{fig:blau_sigma}
\end{figure}

\begin{table}
	\centering
	\caption{Die Messwerte für $\Delta s$ und $\delta s$, sowie die berechneten Werte $\delta\lambda$ für die Aufspaltung der blauen Cadmium $\sigma$-Linien.}
	\label{tab:tabblausigma}
	\sisetup{table-format=1.2}
	\begin{tabular}{S[table-format=3.0]S[table-format=3.0]S[table-format=1.2]}
		\toprule
		{$ \Delta s_{\text{anom}}/ px$} & {$ \delta s_{\text{anom},\sigma}/ px $} & {$  \delta\lambda_{\text{anom},\sigma}/ \si{\pico\meter}$} \\
		\midrule
		356 & 164 & 6.21 \\
		345 & 155 & 6.05 \\
		316 & 146 & 6.23 \\
		301 & 139 & 6.22 \\
		290 & 133 & 6.18 \\
		269 & 128 & 6.41 \\
		260 & 118 & 6.12 \\
		248 & 117 & 6.36 \\
		240 & 110 & 6.18 \\
		218 &  99 & 6.12 \\
		\bottomrule
	\end{tabular}

	\label{tab:blau_sigma}
\end{table}

\subsubsection{Untersuchung der $\pi$-Linien}

Das Magnetfeld wird mit $\SI{18}{\ampere}$ induziert, was mit den Werten der Kalibrierung einem Magnetfeld von $\SI{1.085}{\tesla}$ entspricht. In Abbildung \ref{fig:blau_pi} ist die Aufnahme der Spektrallinien mit und ohne Magnetfeld zu sehen. Die gemessenen Werte für $\Delta s$ und $\delta s$ sind zusammen mit den nach Formel \eqref{eq:deltaLambda} berechneten Werten für $\delta\lambda$ in Tabelle \ref{tab:blau_pi} eingetragen.
Über alle Werte von $\delta\lambda$ gemittelt ergibt sich:
\[
\delta\lambda_.{anom,\pi} = \SI{5.85(5)}{\pico\metre}\text{.}
\]
Mit Formel \eqref{eq:exgj} ergibt sich für die $g_{ij}$:
\[
g_{ij_.{anom,\pi}} = \pm\num{0.502(5)}\text{.}
\]

\begin{figure}
	\centering
	\includegraphics[width=\linewidth-130pt,height=\textheight-130pt,keepaspectratio]{content/images/blau_pi.jpg}
	\caption{Die Aufspaltung der blauen Cadmium $\pi$-Linien.}
	\label{fig:blau_pi}
\end{figure}

\begin{table}
	\centering
	\caption{Die Messwerte für $\Delta s$ und $\delta s$, sowie die berechneten Werte $\delta\lambda$ für die Aufspaltung der blauen Cadmium $\pi$-Linien.}
	\label{tab:tabblaupi}
	\sisetup{table-format=1.2}
	\begin{tabular}{S[table-format=3.0]S[table-format=3.0]S[table-format=1.2]}
		\toprule
		{$ \Delta s_{\text{anom}}/ px$} & {$ \delta s_{\text{anom},\pi}/ px $} & {$  \delta\lambda_{\text{anom},\pi}/ \si{\pico\meter}$} \\
		\midrule
		356 & 161 & 6.09 \\
		345 & 142 & 5.55 \\
		316 & 140 & 5.97 \\
		301 & 130 & 5.82 \\
		290 & 122 & 5.67 \\
		269 & 121 & 6.06 \\
		260 & 112 & 5.80 \\
		248 & 108 & 5.87 \\
		240 & 105 & 5.90 \\
		218 &  94 & 5.81 \\
		\bottomrule
	\end{tabular}

	\label{tab:blau_pi}
\end{table}